\section{Introduction}
Retrieval-augmented generation supplies a language model with passages and asks it to answer from them. Two internal questions matter for such systems: has the model registered that the supplied evidence is now \emph{sufficient}, and has it \emph{integrated} the decisive passage into its answer computation? A reliable pre-generation readout of either would let a system decide whether to retrieve again, detect ignored evidence, or abstain; a causal handle on either would let a system correct the model's behaviour directly. Prior work has shown that residual-stream activations carry signals about context-versus-parametric conflict and context reliance, but these are typically identified on single contexts, and it is rarely tested whether the directions that read such states also control them.

We take a different unit of analysis: the \emph{paired update}. For each multi-hop question we construct nested evidence contexts $C_0\subset C_1\subset\cdots$ and run a fresh, stateless forward pass on each, reading the residual stream at the fixed answer-start position. The difference $\Delta h_{k}=h(C_k)-h(C_{k-1})$ isolates what one added paragraph did to the model's state. Because every decisive increment (the paragraph completing the evidence) is matched with controls that add a length-matched distractor, a duplicate of already-present support, a falsified version of the decisive paragraph, or the answer string embedded in an irrelevant passage from the same penultimate state, a probe trained on these updates cannot succeed by detecting length, answer-string presence or passage relevance alone. Each increment is also labelled with the model's own behavioural transition (wrong$\to$correct, wrong$\to$wrong, correct$\to$wrong, and so on), giving separate supervision for sufficiency, uptake, correction and stability rather than one conflated ``answer-update'' target.

Our contributions are as follows.
\begin{itemize}
  \item \textbf{A matched-increment protocol} for reading evidence-sufficiency and evidence-uptake states from single forward passes, with four-way disjoint evaluation (question, document, answer entity, relation type) and a conversational-revision ablation (\cref{sec:setup,sec:method}).
  \item \textbf{Readout results.} Sufficiency is linearly decodable at $0.95$ AUROC in Qwen2.5-7B, Llama-3.1-8B and Mistral-7B, survives residualising twelve nuisance covariates, orders increment kinds as predicted (decisive $>$ false $>$ distractor $>$ redundant $>$ answer-string control), and transfers without retraining to a conversational protocol ($0.79$), to \twowiki{} ($0.96$) and to 3--4-hop questions ($0.82$). Uptake is decodable at $0.78$ along a direction nearly orthogonal to sufficiency (\cref{sec:readout}).
  \item \textbf{Text baselines.} A fine-tuned DeBERTa that sees the full context recovers sufficiency nearly as well as the probe, but not uptake ($0.56$ vs $0.78$): uptake is the internal signal (\cref{sec:text}).
  \item \textbf{Causal dissociation.} Discriminative probe directions are inert under additive steering (up to 35\% of the residual norm), projection-out and component patching, in all three models. The full anchor residual is causally sufficient, and its effect lives in the orthogonal complement, a high-rank subspace whose composition is model-specific. The mean-difference sufficiency direction is causally active but acts as an abstention gate rather than an accommodation mechanism (\cref{sec:causal}).
\end{itemize}
The picture that emerges is that evidence sufficiency and uptake are cleanly \emph{readable} before generation, that the readout of uptake is not recoverable from text, and that neither readout is a \emph{lever}: the linear probe finds a correlate of the state, while the mechanism that turns evidence into an answer is high-dimensional and content-bearing.
